\section{Introduction}

Language agents are commonly trained to predict which words come next, even when their task is to choose an action according to its consequences. We ask whether an agent can instead predict the latent consequence of each language action and plan with those predicted states. This question separates two capabilities that autoregressive generation usually entangles. A model must represent what has happened, and it must order the actions that remain.

Joint-embedding predictive architectures provide a natural substrate for the first capability \citep{lecun2022path,assran2023ijepa}. An action-conditioned predictor maps the current representation and a language action to a predicted successor representation. Such a model can preserve consequence-relevant information without reconstructing every word. Action-conditioned JEPAs have already supported planning in physical environments \citep{assran2025vjepa2}. Language actions raise a focused question. If the next latent state is predicted accurately, does its geometry support choosing an action?

The answer is not automatic. Predictive learning can identify which histories have the same controlled future, but it need not identify a metric that ranks different futures by progress toward a goal. A bijective re-embedding preserves every latent transition and can still reverse the distances of two successor states to the goal. Prediction gives a state representation. It does not by itself give a planning geometry.

We study this gap in a controlled language-action testbed. Each stylized iGSM episode presents an arithmetic dependency problem. At each step, every learned method receives the same shuffled menu of executable intent phrases. The environment executes the selected phrase and returns its textual consequence. Supplying the current feasible menu is an intentional laboratory intervention. It holds proposal and syntactic validity fixed so that the experiment compares how learning objectives value the same language actions.

We introduce Geometric Advantage Ranking (GAR). GAR uses factual and counterfactual action outcomes to impose local order constraints among candidate consequences. In the main JEPA, a scorer consumes the predicted successor. In a matched direct GAR control, the scorer consumes the history and action without using the predicted successor in its policy path. This terminology separates the supervision from the architectural factorization.

Our empirical goal is explanatory rather than competitive. A predictive GAR JEPA attains $43.3\%$ strict success over five seeds. A token language model attains $42.2\%$, a sentence language model attains $47.1\%$, and a sentence model with an auxiliary latent objective attains $52.1\%$. The JEPA is therefore a viable planner under the shared interface, but it is not the best model. The more informative result is that latent prediction alone attains $18.3\%$. The full method attains $43.3\%$, and removing the GAR terms attains $7.6\%$. Counterfactual ordering is the main source of the gain, while counterfactual successor prediction provides an additional benefit.

The paper develops one connected argument.
\begin{enumerate}
    \item We define the state needed for language-mediated control as a quotient of histories that have identical controlled futures, and show why this state is control sufficient.
    \item We show that exact latent transition prediction does not identify a goal-directed metric without further assumptions. GAR supplies local counterfactual order constraints that are sufficient for strict planning when they recover the optimal action ordering.
    \item We test the mechanism with five-seed objective ablations, gradient-routing controls, matched GAR scorers, and representation diagnostics. The evidence supports GAR as useful ordering supervision. It does not yet establish that successor factorization is uniformly better than direct scoring.
\end{enumerate}

The resulting claim is deliberately narrow: action-conditioned latent prediction can support language-action planning, but predictive accuracy alone leaves the decision geometry underdetermined.
